\begin{abstract}
Masked Diffusion Models (MDMs) offer greater flexibility in decoding order than autoregressive models but require careful planning to achieve high-quality generation. Existing samplers typically adopt greedy heuristics, prioritizing positions with the highest local certainty to decode at each step. Through failure case analysis, we identify a fundamental limitation of this approach: it neglects the downstream impact of current decoding choices on subsequent steps and fails to minimize cumulative uncertainty. In particular, these methods do not fully exploit the non-causal nature of MDMs, which enables evaluating how a decoding decision reshapes token probabilities/uncertainty across all remaining masked positions. To bridge this gap, we propose the \textbf{Info-Gain Sampler}, a principled decoding framework that balances immediate uncertainty with information gain over future masked tokens. 
% We also propose an efficient implementation that ensures practical overhead is minimal. 
Extensive evaluations across diverse architectures and tasks—including reasoning, coding, creative writing, and image generation—demonstrate that the Info-Gain Sampler consistently outperforms existing greedy samplers for MDMs. For instance, it achieves a 3.6\% improvement in average accuracy on reasoning tasks and a 63.1\% win-rate in creative writing. Notably, on reasoning tasks, it reduces cumulative uncertainty from 78.4 to 48.6, outperforming the best greedy baseline by a large margin. We also demonstrate its effectiveness against a concurrent method with a lookahead objective. The code will be available at \url{https://github.com/yks23/Information-Gain-Sampler}.
\end{abstract}